% !TEX root = ../main.tex
% ============================================================================
% Written last, once the numbers of Section 5 are final.
%
% Shape (150 to 200 words): problem (1 sentence), approach (1 to 2), setup (1),
% headline numbers (2 to 3), what they mean (1).
% Rubric gate 1: every claim here points to a number, table, or figure in the body.
% ============================================================================
\begin{abstract}
Monocular depth models such as Depth Anything V2 return a depth value for every pixel and no indication of which values to trust. We train a small convolutional head that reads the frozen depth model's decoder features, its predicted depth, the temporal spread of depth across neighboring frames, and two radar channels, and outputs per cell the probability that the depth error exceeds a fixed threshold. Labels come from lidar projected into the camera image, so no hand annotation is needed. On nuScenes, we train 18 heads that differ in their inputs, radar settings, label threshold, and training data, on a cache of all 204,894 trainval camera images.\footnote{\raggedright Code: \url{https://github.com/kgqiao/monocular-depth-confidence-estimation}} \todo{Two or three sentences with the headline numbers once the scores are final. (1) AUSE and AURG of the best head that needs no lidar at inference, against flip disagreement and the oracle. (2) The gap between day, rain, and night. (3) What the input ablation shows: decoder features, temporal spread, radar. Then one closing sentence on what this means for a consumer of the depth.}
\end{abstract}
